\section{Data and event construction}

\subsection{Primary public-data corpus}

The primary corpus comprises EPA AirData daily archives for parameter 88101,
PM2.5 local conditions, for 2019--2025
\citep{epa_obtaining_aqs,epa_airdata_formats}. The canonical signal is the
\texttt{Arithmetic Mean} for \texttt{24-HR BLK AVG} observations. The
predeclared cleaning rule retains nonempty Method Code records with observation
percent at least 75 and an Event Type other than \texttt{Excluded}. It prevents
alternative daily summaries for the same monitor-day from being treated as
independent observations.

\begin{table}[tbp]
\centering
\caption{Frozen v0.3.2 data and audit inventory.}
\label{tab:data-summary}
\small
\begin{tabular}{lr}
\toprule
Quantity & Count \\
\midrule
Canonical daily PM2.5 records & 2,424,793 \\
Monitor time series & 1,689 \\
Stable pseudo-anchors from Method Code regimes & 124 \\
Stable pseudo-anchors from all-monitor regimes & 22 \\
Persistent Method Code anchors & 563 \\
Anchors with at least one distinct physical donor & 394 \\
Anchors with at least three distinct physical donors & 238 \\
Complete common-method comparisons & 228 \\
Donor-insufficient anchors & 325 \\
Estimator input-window failures & 10 \\
\bottomrule
\end{tabular}
\par\smallskip
\begin{minipage}{0.94\linewidth}
\footnotesize\textit{Note.} The physical donor identifier is State Code + County Code + Site Number; at most one POC is retained per donor site.
\end{minipage}
\end{table}


The source files, archive size, checksum, selected CSV member, and retained row
count are preserved in the frozen data manifest. The daily data are public, but
raw downloads and API responses are not included in this manuscript repository
or its paper assets. That separation keeps the report reproducible without
placing raw files or local credentials under version control.

\subsection{Metadata anchors and distinct physical donors}

An anchor is a persistent reported Method Code transition with at least 60
calendar days on each side, at least 45 retained observations in each adjacent
method run, and no more than a seven-day transition gap. An eligible geographic
donor is at a different physical site, within the predeclared distance limit,
method-stable in the target window, sufficiently paired in the pre-event
period, and sufficiently correlated before the anchor.

Physical donor identity is State Code + County Code + Site Number. A target's
candidate donor list can contain multiple POCs for one physical site, but the
ranked donor rule retains at most one: highest pre-event correlation, then
shorter distance, then more paired days, then POC as a deterministic
tie-breaker. This rule remediates an earlier superseded inventory in which
multiple POCs could be treated as independent geographic donors. Same-site
alternate POCs are excluded from geographic donor construction and reserved for
contextual consistency analysis.

\subsection{Audit disposition}

Every anchor is retained in the event audit. A complete comparison means that
the required common cross-site inputs exist; donor insufficiency and
input-window failure are explicit machine-readable outcomes rather than omitted
rows. This complete accounting is necessary because a result on only the
well-matched subset would overstate the scope of cross-site inference.
